% 2011 年陕西卷理科第 5 题：正方体挖去圆锥的三视图。
% 来源 PDF：普通高考/2011/2011陕西理.pdf，SHA-256 04dab6de0e3ee6223a1cbc126f7f83ae354d2b31201a14db1f31c8477d8a3bed。
% 原页：第 1 页第 1 栏下部，书页码 693。外形 2×2×2，俯视图内切圆半径 1。
\begin{tikzpicture}[semithick, font=\small, scale=0.85]
  % 主视图
  \draw[thick] (0,0) rectangle (2,2);
  \draw[dashed] (0,2)--(1,0.35)--(2,2);
  \node at (1,-0.42) {主视图};
  \draw[<->] (0,2.42)--(2,2.42) node[midway, above] {$2$};
  % 主视图宽度标注
  % 左视图
  \begin{scope}[xshift=3.6cm]
    \draw[thick] (0,0) rectangle (2,2);
    \draw[dashed] (0,2)--(1,0.35)--(2,2);
    \node at (1,-0.42) {左视图};
    \draw[<->] (2.4,0)--(2.4,2) node[midway, right] {$2$};
  \end{scope}
  % 俯视图
  \begin{scope}[yshift=-3.6cm]
    \draw[thick] (0,0) rectangle (2,2);
    \draw[thick] (1,1) circle (1);
    \node at (1,-0.42) {俯视图};
  \end{scope}
\end{tikzpicture}
